% SEGMENT OLP-0540-S01
% Part: set-theory
% Chapter: z
% Section: union

\documentclass[../../../include/open-logic-section]{subfiles}

\begin{document}

\olfileid{sth}{z}{union}
\olsection{சேர்ப்பு}

\olref[sth][z][sep]{prop:intersectionsexist} நமக்கு வெட்டுகளைத் தந்தது.
ஆனால் தன்னிச்சையான சேர்ப்புகள் இருக்க வேண்டுமெனில், மேலும் ஓர் அடிகோளை
வகுக்க வேண்டும்:

\begin{axiom}[சேர்ப்பு] எந்தக் கணம் $A$ க்கும், பின்வரும் கணம் உள்ளது:
$\bigcup A = \Setabs{x}{(\exists b \in A) x \in b}$.
\[
	\forall A \exists U \forall x(x \in U \liff (\exists b \in A)x \in b)
\]
\end{axiom}

திரள்வுறும் படிநிலைமுறைக் கருத்தாக்கமும் இந்த அடிகோளை நியாயப்படுத்துகிறது.
$A$ ஒரு கணம் என்க; ஆகவே $A$ ஏதோ ஒரு படிநிலை $S$ இல் உருவாகிறது
(\stageshier{} இன்படி). $A$ இன் ஒவ்வோர் உறுப்பும் $S$ க்கு \emph{முன்பு}
உருவானது (\stagesacc{} இன்படி); அதேவிதமாகப் பகுத்தறிந்தால், $A$ இன்
ஒவ்வோர் உறுப்பினுடைய ஒவ்வோர் உறுப்பும் $S$ க்கு முன்பு உருவானது. ஆகவே
\emph{அந்தக்} கணங்கள் அனைத்தும் $S$ க்கு முன்பு கிடைக்கின்றன; அவற்றை
$S$ இல் ஒரு கணமாக உருவாக்கலாம். அந்தக் கணமே $\bigcup A$.

\end{document}
